\section{Endocitosis}

Es un tipo de transporte activo que mueve partículas hacia una célula. Se caracteriza porque su mecanismo consiste en la invaginación de la membrana plasmática para formar un "bolsillo" alrededor de la partícula diana.

\textbf{Partículas transportadas:}

\begin{itemize}
    \item Moléculas grandes
    \item Partes de una célula
    \item Células enteras
\end{itemize}

\section{Variaciones de endocitosis}

\subsection{Fagocitosis}
Proceso de "comer" células o partículas grandes. Es llevado a cabo por células especializadas como macrófagos y neutrófilos.

\subsection{Pinocitosis}
Variante donde se consumen partículas pequeñas junto con agua del líquido extracelular. No necesita de lisosomas. Es un proceso constitutivo en muchas células.

\subsection{Endocitosis mediada por receptores}
Es un tipo de endocitosis específica donde las partículas se unen a receptores en la superficie celular antes de ser internalizadas.

\section{Mecanismo de endocitosis}

\begin{enumerate}
    \item Comienza con la proteína clatrina adhiriéndose a la membrana.
    \item Esta porción se extiende para rodear la partícula y crea una vesícula.
        \begin{itemize}
            \item Endosoma: Vesícula en formación con contenido extracelular
            \item Fagosoma: Vesícula que contiene partículas grandes
        \end{itemize}
    \item Se acidifica el ambiente para "digerir" la partícula y se ponen en acción los lisosomas.
        \begin{itemize}
            \item Enzimas: Proteasas, lipasas y nucleasas
            \item Moléculas antimicrobianas
        \end{itemize}
    \item Se libera el material digerido (exocitosis).
\end{enumerate}

\section{Exocitosis}

Es el proceso inverso a la endocitosis. Se liberan vesículas con material hacia el exterior de la célula. Es importante para:
\begin{itemize}
    \item Secreción de hormonas
    \item Liberación de neurotransmisores
    \item Eliminación de desechos
\end{itemize}